\clearpage
\thispagestyle{empty}

\begin{flushright}
  {\fontsize{28pt}{28pt}\selectfont\bfseries Agradecimientos\par}
\end{flushright}

\vspace{0.3cm}

\begingroup
\leftskip=0.10\textwidth
\parindent=0pt
\fontsize{10pt}{12.5pt}\itshape\selectfont

Quiero comenzar agradeciendo a mi familia por haber estado siempre ahí y por
haberme permitido estudiar en la universidad. Su apoyo constante ha sido la base
para afrontar este esfuerzo. 

A mi compañera Alba, por el trabajo conjunto y por el apoyo durante el primer
semestre.

A mi compañero Guillermo, por todo el apoyo durante estos cuatro años. Gracias a
él he participado en actividades y montado proyectos que, sin su impulso, no hubieran
salido adelante.

A mis amigos Alejandro y Alberto, compañeros de \texttt{refinement} continuo. Gracias
a ellos formamos una trifecta con la que navegamos las novedades tecnológicas mes tras
mes, perfeccionando nuestro estilo, nuestros procedimientos y nuestra forma de mejorar.
Porque no sumamos tres, multiplicamos ocho. Gracias por todas las largas horas
peleándonos con Gemini CLI u Opencode para conseguir que colocase una imagen en
LaTeX exactamente como queríamos; por todos los intentos con herramientas chinas
para convertir imágenes en texto seleccionable; por los experimentos creando
agentes personalizados; por el trabajo con conexiones MCP y con CLI para ampliar las
capacidades de la inteligencia artificial; y por todos los intentos de lograr aquello que
nos decían que era imposible.

A mis tutores del TFG, Joserra y Guzmán, por permitirme desarrollar el trabajo
que deseaba y por darme margen para hacerlo con herramientas de vanguardia,
aprovechando todos los conocimientos que había ido adquiriendo. Su confianza
hizo posible que este proyecto fuera muy ambicioso, coherente con mi manera de
avanzar, experimentar y desarrollar. 

A mi profesor Andrés Merino, de Meteorología Aeronáutica, por no solo permitir,
sino motivar, que en la parte práctica de su asignatura se utilizase IA generativa
en programación para alcanzar resultados mucho más avanzados. Todo ello se apoyaba en
la misma teoría que habría sustentado un trabajo convencional, pero gracias a
las herramientas pudimos construir simuladores más complejos, montar estaciones
meteorológicas y crear aplicaciones desplegadas en la nube dignas de ser presentadas
en concursos. Gracias por permitirme desarrollar mis capacidades en un entorno
real de ingeniería y por darme espacio para perfeccionar mis técnicas, mis
procedimientos y mis conocimientos sobre cómo gestionar estas herramientas.

Y, finalmente, a todas las personas que me escuchan, que me apoyan y que están
ahí en los momentos más complicados. A todas las personas que me dijeron
``pruébalo'', ``inténtalo'', ``tú puedes''. Gracias, porque demostramos
que sí era posible.

\par
\endgroup

\clearpage
\thispagestyle{empty}
\null\vfill
\begin{center}
  \textit{This page intentionally left blank.}
\end{center}
\vfill\null
